\section{An integrated growing central band for the physical record}
\label{sec:growing-major-arc}
The estimate on compact rescaled frequency sets can be strengthened to
an integrated estimate on an expanding band. The distinction matters
for local inversion: convergence after integration controls the inverse
transform uniformly in the observation variable. We prove this estimate
for the actual return record, including bounded-variation initial
insertions. No induced spectral realization is used.

An admissible initial insertion is a complex-valued function $a_R$ on
$M$, vanishing almost everywhere outside $Y_R^*$, such that
\[
 \|a_R\|_\infty+\|a_R\|_{\mathrm{BV}}\le B.
\]
Its integral is denoted by $\alpha_R$. The original probability has
$a_R=s_R$ and $\alpha_R=1$. Write
\begin{equation}\label{eq:inserted-central-characteristic}
 C^a_{n,R}(v)=\int_{Y_R^*}a_R(x)
 \exp\left(\frac{i v\cdot(J_{n,R}(x)-n\bar G_R)}{\sqrt n}\right)
 \dd\nu(x),\qquad v\in\R^4.
\end{equation}
Here $v$ is the rescaled frequency; the physical Fourier frequency is
$\omega=v/\sqrt n$. Constants below depend on $B$ but not on the
insertion, the radius, or $n$. No continuity of the insertion in $R$
is required.

\subsection{A frequency-dependent error bound}
\begin{lemma}[Collision estimate with explicit frequency dependence]
\label{lem:explicit-frequency-error}
Put $\ell_\delta=1+|\log\delta|$. There are $c,C>0$ and $\rho<1$
such that, for $m\ge2$, $0<\delta<1/2$, and real $t$ satisfying
\[
 |t|\le c\sqrt m\,\delta^2,\qquad
 \delta^{-6}|t|^3m^{-1/2}\le c,
\]
every complex function $a$ with
$\|a\|_\infty+\|a\|_{\mathrm{BV}}\le B$ satisfies
\begin{align}
 &\left|\int a\exp(it\cdot S_{m,R}h_R/\sqrt m)\dd\nu
       -\left(\int a\dd\nu\right)e^{-t^{\mathsf T}\Gamma_Rt/2}\right|
       \notag\\[-2mm]
 &\quad\le C\left\{
 \delta+|t|\sqrt{\delta\ell_\delta}+|t|^2\delta\ell_\delta
 +\frac{|t|\delta^{-4}+|t|^3\delta^{-6}}{\sqrt m}
 +\delta^{-2}\rho^m\right\}.\label{eq:explicit-frequency-error}
\end{align}
\end{lemma}
\begin{proof}
Set $a_\delta=\mathsf S_\delta a$ and $z=t/\sqrt m$.
The operator in Lemma~\ref{lem:smooth-collision-expansion} gives
\[
 \int a_\delta e^{it\cdot S_{m,R}h_{R,\delta}/\sqrt m}\dd\nu
 =\lambda_{R,\delta}(z)^m
   \ell_j\bigl(\Pi_{R,\delta}(z)a_\delta\nu\bigr)
       +O(\delta^{-2}\rho^m).
\]
The embedding bound in Lemma~\ref{lem:density-norm-details} below
makes the initial-vector factor explicit. The principal amplitude
at zero is $\int a\dd\nu$, and its change is at most
$C|t|\delta^{-4}/\sqrt m$. Moreover
\[
 m\log\lambda_{R,\delta}(t/\sqrt m)
 =-\tfrac12t^{\mathsf T}\Gamma_{R,\delta}t
       +O(\delta^{-6}|t|^3/\sqrt m).
\]
The stated smallness conditions put $z$ in the analytic ball and bound
the last remainder. Positive semidefiniteness of
$\Gamma_{R,\delta}$ therefore bounds the principal power by a fixed
constant. The inequality $|e^w-1|\le |w|e^{|w|}$ proves the two
spectral error terms in \eqref{eq:explicit-frequency-error}.

Mean-preserving smoothing changes the initial insertion by $O(\delta)$
in $L^1$. If $r=h_R-h_{R,\delta}$, then
\[
 \int |a_\delta|\,|S_{m,R}r|^2\dd\nu
       \le C m\delta\ell_\delta,
\]
by \eqref{eq:small-bv-variance} and $|a_\delta|\le B$.
Cauchy--Schwarz bounds the characteristic error from removing this
smoothing by $C|t|\sqrt{\delta\ell_\delta}$. This argument uses
absolute values of the insertion and so applies to signed and complex
insertions as well. Finally, the segment joining the two covariance
matrices consists of positive semidefinite matrices. Differentiation
along that segment and \eqref{eq:collision-covariance-limit} give
\[
 \left|e^{-t^{\mathsf T}\Gamma_{R,\delta}t/2}
             -e^{-t^{\mathsf T}\Gamma_Rt/2}\right|
       \le C|t|^2\delta\ell_\delta.
\]
Combining these estimates proves the lemma.
\end{proof}

\begin{proposition}[Stopped estimate on finite frequency balls]
\label{prop:stopped-band-error}
For $U\ge1$, let
\[
 \mathcal E_n(U)=\sup_{R,a_R}
   \int_{|v|\le2U}
    |C^a_{n,R}(v)-\alpha_Re^{-v^{\mathsf T}D_Rv/2}|\dd v.
\]
If $U\le c\sqrt n\,\delta^2$ and
$\delta^{-6}U^3/\sqrt n\le c$, with $c$ sufficiently small, then
\begin{align}
 \mathcal E_n(U)\le C\bigl\{
 &\delta U^4+\sqrt{\delta\ell_\delta}\,U^5
       +\delta\ell_\delta U^6
       +n^{-1/2}\delta^{-4}U^5+n^{-1/2}\delta^{-6}U^7
       \notag\\
 &+\delta^{-2}\rho_0^nU^4
       +n^{-1/5}U^4+n^{-1/5}\log(2+n)U^5+n^{-1}U^6
                          \bigr\},\label{eq:integrated-band-error}
\end{align}
where $\rho_0<1$. The supremum is over the admissible insertions.
\end{proposition}
\begin{proof}
Put $m_n=\lfloor n/c_*\rfloor$ and $b_n=\lceil n^{3/5}\rceil$.
The exact compensation identity \eqref{eq:exact-stopped-compensation}
permits replacement of the centered record by $S_{N_{n,R},R}h_R$.
The exceptional set $|N_{n,R}-m_n|>b_n$ has $\nu_R^*$-probability
$O(n/b_n^2)$ by Lemma~\ref{lem:stopped-compensation}; rounding changes
only the constant. Its contribution with insertion $a_R$ is bounded
by the same order, since $|a_R|\nu\le Bc_*\nu_R^*$ on the section.
On its complement use the forward and backward deterministic-window
maxima from Lemma~\ref{lem:dyadic-collision-maximal}. The elementary
exponential difference bound and Cauchy--Schwarz give
\begin{align}
 &\left|C^a_{n,R}(v)
       -\int a_R e^{iv\cdot S_{m_n,R}h_R/\sqrt n}\dd\nu\right|
       \notag\\
 &\hspace{12mm}\le C\left\{n^{-1/5}
                  +|v|n^{-1/5}\log(2+n)\right\}.
                 \label{eq:weighted-stopping-error}
\end{align}
No second moment of an unbounded stopped sum on the exceptional set
is used. In particular the estimate is valid for complex insertions.

Apply Lemma~\ref{lem:explicit-frequency-error} at $m_n$ with
$t=v\sqrt{m_n/n}$. These frequencies satisfy its restrictions after
changing the uniform constant $c$. The covariance of the deterministic
sum is $(m_n/n)\Gamma_R$, which differs from $D_R=c_*^{-1}\Gamma_R$
by $O(n^{-1})$. Its Gaussian characteristic function consequently
changes by at most $C|v|^2/n$. Integrating these bounds and
\eqref{eq:weighted-stopping-error} uses only
$\int_{|v|\le2U}|v|^j\dd v=C_jU^{4+j}$, for $j\ge0$.
This proves \eqref{eq:integrated-band-error}. Finitely many small values
of $n$ can be absorbed by increasing the constants wherever the stated
smallness conditions hold.
\end{proof}

\subsection{A polynomially expanding major arc}
\begin{theorem}[Integrated major arc of the actual return record]
\label{thm:growing-major-arc}
Let
\[
 0<\varepsilon<\frac1{134},\qquad
 10\varepsilon<\theta<\frac{1/2-7\varepsilon}{6},\qquad
 r=\min\left\{\frac\theta2-5\varepsilon,
              \frac12-6\theta-7\varepsilon,
              \frac15-5\varepsilon\right\}.
\]
Then $r>0$ and
\begin{equation}\label{eq:general-growing-arc}
 \sup_{R,a_R}\int_{|v|\le2n^\varepsilon}
 |C^a_{n,R}(v)-\alpha_Re^{-v^{\mathsf T}D_Rv/2}|\dd v
       \le C n^{-r}\log(2+n).
\end{equation}
In particular, with $\varepsilon=1/200$ and $\theta=1/14$,
\begin{equation}\label{eq:explicit-growing-arc}
 \sup_{R,a_R}\int_{|v|\le2n^{1/200}}
 |C^a_{n,R}(v)-\alpha_Re^{-v^{\mathsf T}D_Rv/2}|\dd v
       \le C n^{-3/280}\sqrt{\log(2+n)}.
\end{equation}
These estimates hold for $n\ge2$. They require only the positive
semidefinite covariance already established, not positive definiteness.
\end{theorem}
\begin{proof}
Take $U=n^\varepsilon$ and $\delta=\tfrac14n^{-\theta}$ in
Proposition~\ref{prop:stopped-band-error}. Its analytic restrictions hold
for all sufficiently large $n$: in particular
$3\varepsilon+6\theta<1/2$ follows from
$7\varepsilon+6\theta<1/2$, and
$\varepsilon+2\theta<1/2$ follows as well. Constants in the radius
of the analytic ball do not affect these strict inequalities.

The three potentially slow terms in \eqref{eq:integrated-band-error}
are, respectively,
\[
 n^{-(\theta/2-5\varepsilon)}\sqrt{\log(2+n)},\qquad
 n^{-(1/2-6\theta-7\varepsilon)},\qquad
 n^{-(1/5-5\varepsilon)}\log(2+n).
\]
All other polynomial terms decay at least as fast as one of these;
this follows by subtracting their exponents, using $\theta>10\varepsilon$.
For the rounding term also note that
$1-6\varepsilon>1/5-5\varepsilon$ in the indicated range.
The exponentially decaying term is smaller than every fixed negative
power. The interval for $\theta$ is nonempty precisely when
$67\varepsilon<1/2$, proving \eqref{eq:general-growing-arc}.

At the displayed explicit choice the three margins are
\[
 \frac3{280},\qquad \frac{51}{1400},\qquad \frac7{40}.
\]
Only the first is minimal, and that term carries a square root of the
logarithm. The strict excess of the other margins absorbs their
logarithmic factors. Increasing the constant handles the finite initial
range, since both characteristic functions are bounded uniformly by
a constant depending on $B$. This proves \eqref{eq:explicit-growing-arc}.
\end{proof}

\subsection{Insertion into exact raw inversion}
Define the finite complex measure $\mu^a_{n,R}$ by pushing forward
$a_R\nu$ under $J_{n,R}$, and write $\Phi^a_{n,R}$ for its transform.
Take $a_n=n^{1/200}$ in the cutoff and kernel of
\eqref{eq:localized-kernel}, and use $m_R=\bar G_R$ in the central
windows. The preceding theorem supplies a proved replacement for the
spectral small-frequency hypothesis in the following inversion route.

\begin{corollary}[Raw inversion from the proved central band]
\label{cor:raw-from-growing-arc}
Suppose, in addition, that $D_R\ge d_0 I_4$ for some $d_0>0$.
For an admissible family of insertions, let $E^a_{n,R}$ be finite complex
measures with mixed densities $e^a_{n,R}$ and put
$H^a_{n,R}=\Phi^a_{n,R}-\widehat E^a_{n,R}$. Assume that
$H^a_{n,R}\in L^1(\T^3\times\R)$ and, uniformly in the family,
\begin{equation}\label{eq:new-complementary-tail}
 \int |1-\chi_n(\omega)|\,|H^a_{n,R}(\omega)|\dd\omega=o(n^{-2}).
\end{equation}
Assume also the two local edge conditions
\eqref{eq:local-edge-errors} with $e,E$ replaced by $e^a,E^a$.
Then $\mu^a_{n,R}$ has a mixed density $p^a_{n,R}$, and for each fixed
$K$,
\begin{align}
 \mathop{\rm ess\,sup}_{|z-n\bar G_R|\le K\sqrt n}
 \left|n^2p^a_{n,R}(z)
 -\frac{\alpha_R}{(2\pi)^2\sqrt{\det D_R}}
   \exp\left(-\tfrac12\xi^{\mathsf T}D_R^{-1}\xi\right)\right|
 &\longrightarrow0,
 \quad \xi=\frac{z-n\bar G_R}{\sqrt n},
 \label{eq:inserted-raw-conclusion}
\end{align}
uniformly in $R$ and the admissible family. No fixed-neighborhood
spectral expansion of an induced operator is assumed in this corollary.
\end{corollary}
\begin{proof}
Fourier inversion for $H^a$ identifies a continuous density for
$\mu^a-E^a$. Add $e^a$ and use the exact decomposition
\eqref{eq:exact-local-decomposition}, which is linear and remains valid
for complex measures. After $v=\sqrt n\,\omega$, the difference between
$n^2K_n*\mu^a$ and the inverse Gaussian integral with cutoff
$\chi(v/n^{1/200})$ is bounded, independently of $z$, by
$(2\pi)^{-4}\|\chi\|_\infty$ times the integral in
\eqref{eq:explicit-growing-arc}. Positive definiteness bounds the
removed Gaussian tail by
$C\int_{|v|>n^{1/200}}e^{-d_0|v|^2/2}\dd v$, which tends to zero
uniformly. This proves the main term in
\eqref{eq:inserted-raw-conclusion}. The other two terms of the exact
decomposition vanish on the stated scale by the assumed complementary
tail and local edge conditions. These steps prove the raw density
conclusion, not a conclusion only for a convolved observation.
\end{proof}

\begin{remark}[Location and scope of the remaining estimates]
\label{rem:growing-arc-boundary}
The physical radius of the proved central band is
$n^{-1/2+1/200}$. It is smaller than $n^{-2/5}$. Thus
\eqref{eq:new-complementary-tail} is not a consequence of the outer-tail
hypothesis in Theorem~\ref{thm:LLT}: the intervening annulus also needs
an estimate. The present theorem removes the need for an assumed induced
spectral expansion on the central band only. It does not assert a
resolvent estimate on its complement, uniform positive definiteness,
or the full physical branch decomposition. Both inversion routes are
retained, with their distinct frequency boundaries.

The insertion class is an initial-state class with a uniform norm on
the collision cylinder. Terminal insertions, multiple-time insertions,
and exact-event indicators need not belong to it. In particular
\eqref{eq:explicit-growing-arc} does not justify a change of exact
conditioning event or provide the weighted raw tail required for such
a change. The count truncation and the initial-coordinate variation
bounds continue to concern different objects from the derivative sums
of the raw residual densities.
\end{remark}
